% ================================================================
% SECTION 10: FLUX JUMP MODULE
% ================================================================
\section{Flux Jump Module}
\label{sec:flux_jump}

\begin{center}
\code{bionetflux.core.flux\_jump} \quad(327 lines)
\end{center}

The module contains a single public function and an embedded self-test.

% ------------------------------------------------------------------
\subsection{domain\_flux\_jump}

\begin{lstlisting}[language=Python, caption={Flux jump signature}]
def domain_flux_jump(
    trace_solution: np.ndarray,       # neq*(N+1) x 1
    forcing_term: np.ndarray,         # (2*neq) x N
    problem,                          # unused (backward compat.)
    discretization,                   # unused (backward compat.)
    static_condensation               # StaticCondensationBase
) -> Tuple[np.ndarray, np.ndarray, np.ndarray, np.ndarray]:
    """Returns (U, J, F, JF)."""
\end{lstlisting}

\noindent The arguments \code{problem} and \code{discretization} are retained in the signature for backward compatibility but are not used internally; all information is obtained from \code{static\_condensation}.

\subsubsection{Return Values}

\begin{longtable}{p{2.5cm}p{4cm}p{7.5cm}}
\toprule
\textbf{Name} & \textbf{Shape} & \textbf{Description} \\
\midrule
\code{U} & $(2 n_\mathrm{eq}) \times N$ & Bulk polynomial coefficients reconstructed from the trace \\
\code{J} & $(d_f) \times N$ or \code{None} & Flux polynomial coefficients, where $d_f = \sum_k (\mathrm{flux\_order}_k + 1)$ \\
\code{F} & $n_\mathrm{eq}(N{+}1) \times 1$ & Global residual: accumulated flux-trace contributions \\
\code{JF} & $n_\mathrm{eq}(N{+}1) \times n_\mathrm{eq}(N{+}1)$ & Global Jacobian w.r.t.\ trace unknowns \\
\bottomrule
\end{longtable}

\subsubsection{Algorithm}

For each element $k = 0, \ldots, N{-}1$:
\begin{enumerate}
  \item Build a local index array of length $2 n_\mathrm{eq}$, pairing the left and right trace node of each equation.
  \item Extract the local trace column and forcing column.
  \item Call \code{static\_condensation.static\_condensation(local\_trace, gk)}.
  \item Store bulk coefficients in $U[:,k]$ and flux coefficients in $J[:,k]$.
  \item Scatter the returned \code{flux\_trace} into $F$ and the local \code{jacobian} into $JF$ (both via additive assembly using \code{np.add.at}).
\end{enumerate}

\subsubsection{Built-in Self-Test}

\code{test\_domain\_flux\_jump(verbose=True)} creates mock objects (\code{MockProblem}, \code{MockDiscretization}, \code{MockStaticCondensation}) and runs three parametric test cases plus two edge-case tests (zero trace, zero forcing).  The function returns \code{True} if all checks pass.
